% sections/abstract.tex --- Abstract
\begin{abstract}
Machine-learning research on accounting-anomaly detection has produced
a steady stream of methodological advances over the past decade---deep
autoencoders for journal-entry outliers~\citep{schreyer2017detection},
adversarial detectors for accounting deepfakes~\citep{schreyer2019adversarial},
graph neural networks for transaction-graph fraud, and a wide variety of
tabular baselines.  Yet the field has no shared, reproducible benchmark.
Every published result is reported on a private dataset---typically a
single client's general ledger held under audit-confidentiality
constraints---which makes cross-paper comparison impossible and slows
the accumulation of knowledge that public benchmarks
\citep{deng2009imagenet,krizhevsky2009cifar,hu2020ogb} have driven in
neighbouring fields.

We close this gap with \bench{}, an open, fully-provenanced reference
benchmark for accounting-network anomaly detection.  \bench{} is built on
\system{}~\citep{ivertowski2026datasynth}, a Rust framework that
generates end-to-end coherent synthetic enterprise datasets where every
node, edge, and property derives from a known generative process.  The
empirical foundation of the benchmark is an internal study of $364$
million journal-entry postings drawn from $155$ real-world general
ledgers~\citep{ivertowski2024apn}, from which we recover the line-item
distribution and structural network properties that \system{}
implements.

The benchmark contributes four reusable artefacts.  \emph{First}, a
$13$-column flat edge schema (\texttt{graphs/je\_network.csv}) that
materialises five mathematically-equivalent network-construction
methods (Methods~A--E) over the journal-entry table, with
deterministic UUID v5 edge identifiers and a \texttt{predecessor\_line\_id}
field that traces booking chains across documents.
\emph{Second}, a labelled anomaly-injection layer covering more than
$130$ subtypes---spanning fraud, error, process-control, statistical,
and relational categories---with every flagged edge linked back to a
ground-truth event class, enabling both binary and multi-class tasks
without re-labelling.
\emph{Third}, three reproducible evaluation splits---temporal (last
three fiscal months held out), entity (ten percent of subsidiaries held
out), and fraud-typology (typology classes held out)---together with a
leakage audit so that any future submission can be checked against the
same data partition.
\emph{Fourth}, a baseline suite spanning graph neural networks (GCN,
GraphSAGE, GAT, GIN, RGCN), tabular gradient boosting (XGBoost), and
classical anomaly detectors (Isolation Forest, Local Outlier Factor),
together with a simple rule-based reference, all trained from a single
declarative configuration.

We report initial numbers on the released configuration.  Strong GNN
baselines reach competitive area-under-precision-recall on the binary
task, but cross-typology generalisation under the typology-hold-out
split remains substantially harder---suggesting that structural
representations alone do not fully transfer across fraud schemes.
The full benchmark---data-generation configuration, generated dataset,
edge-list export, splits, baselines, and weights---is released under a
permissive licence at the \system{} repository
\citep{datasynth-repo}.

\bench{} does not replace empirical work on real client data.  Its role
is complementary: a reference frame against which methods can be
trained, ablated, and compared without the confidentiality concerns
that block public release of the underlying ledgers.

\medskip
\noindent\textbf{Keywords:} accounting networks, journal entries,
synthetic benchmarks, graph neural networks, anomaly detection, fraud
detection, audit analytics, reproducibility, ISO~21378, reference
knowledge graphs
\end{abstract}
